\documentclass[11pt,twoside]{article}
% Shared formatting for all notes in this repository.
% Adapted from the UC Berkeley EECS CS267 lecture-note template:
% https://www.cs.cmu.edu/~ggordon/10725-F12/template.tex
% Personal-note adaptation: boxed title/author/date; ordinary numbering;
% existing natbib citations and mathematical environments are retained.
% Load after \documentclass; compile from the repository root.
\usepackage[T1]{fontenc}
\usepackage[utf8]{inputenc}
\usepackage{amsmath,amsfonts,amssymb,mathrsfs,amsthm}
\usepackage{graphicx,wrapfig}
\usepackage[inline,shortlabels]{enumitem}
\usepackage{centernot}
\usepackage[numbers]{natbib}
\usepackage{adjustbox,booktabs,arydshln}
\usepackage{caption,subcaption}
\usepackage{xcolor,tikz,pifont}
\usetikzlibrary{decorations.pathreplacing,angles,quotes,shapes.geometric,
  arrows.meta,positioning,calc,fit,backgrounds,shadows}
\usepackage[ruled,vlined]{algorithm2e}
\usepackage{xr}
\usepackage{hyperref}
\hypersetup{colorlinks=true,linkcolor=blue,citecolor=blue,urlcolor=blue}

% Berkeley EECS lecture-note page layout and paragraph spacing.
\setlength{\oddsidemargin}{0.25in}
\setlength{\evensidemargin}{-0.25in}
\setlength{\topmargin}{-0.6in}
\setlength{\textwidth}{6.5in}
\setlength{\textheight}{8.5in}
\setlength{\headsep}{0.75in}
\setlength{\parindent}{0in}
\setlength{\parskip}{0.1in}

% A wrapping box accommodates long research-note titles.
\makeatletter
\renewcommand{\maketitle}{%
  \pagestyle{myheadings}%
  \thispagestyle{plain}%
  \begin{center}
    \fbox{%
      \begin{minipage}{\dimexpr\textwidth-2\fboxsep-2\fboxrule\relax}
        \vspace{2mm}
        \textbf{Research Notes}\hfill\@date\par
        \vspace{2mm}
        {\centering\Large\bfseries\@title\par}
        \vspace{2mm}
        {\itshape Author: \@author\par}
        \vspace{2mm}
      \end{minipage}%
    }
  \end{center}
  \markboth{Research Notes}{Research Notes}%
  {\footnotesize\itshape
    Style adapted from the UC Berkeley EECS lecture-note template.\par}
  \vspace{4mm}%
}
\makeatother

\newcommand{\bR}{{\boldsymbol R}}
\newcommand{\bX}{{\boldsymbol X}}
\newcommand{\bY}{{\boldsymbol Y}}
\newcommand{\by}{{\boldsymbol y}}
\newcommand{\bx}{{\boldsymbol x}} 
\newcommand{\bz}{{\boldsymbol z}}
\newcommand{\bZ}{{\boldsymbol Z}}
\newcommand{\bH}{{\boldsymbol H}}
\newcommand{\bb}{{\boldsymbol b}}
\newcommand{\be}{{\boldsymbol e}}
\newcommand{\bbf}{{\boldsymbol f}}
\newcommand{\bg}{{\boldsymbol g}}
\newcommand{\bm}{{\boldsymbol m}}
\newcommand{\bw}{{\boldsymbol w}}
\newcommand{\bW}{{\boldsymbol W}}
\newcommand{\bv}{{\boldsymbol v}}
\newcommand{\bt}{{\boldsymbol t}}
\newcommand{\bA}{{\boldsymbol A}}
\newcommand{\ba}{{\boldsymbol a}}
\newcommand{\bD}{{\boldsymbol D}}
\newcommand{\bE}{{\boldsymbol E}}
\newcommand{\bi}{{\boldsymbol i}}
\newcommand{\bc}{{\boldsymbol c}}
\newcommand{\bC}{{\boldsymbol C}}
\newcommand{\bk}{{\boldsymbol k}}
\newcommand{\bK}{{\boldsymbol K}}
\newcommand{\bP}{{\boldsymbol P}}
\newcommand{\mK}{\mathcal{K}}
\newcommand{\bV}{{\boldsymbol V}}
\newcommand{\bG}{{\boldsymbol G}}
\newcommand{\mG}{\mathcal{G}}
\newcommand{\mH}{\mathcal{H}}
\newcommand{\mE}{\mathbb{E}}
\newcommand{\mX}{\mathcal{X}}
\newcommand{\mM}{\mathcal{M}}
\newcommand{\bU}{{\boldsymbol U}}
\newcommand{\bu}{{\boldsymbol u}}
\newcommand{\bI}{{\boldsymbol I}}
\newcommand{\bJ}{{\boldsymbol J}}
\newcommand{\bs}{{\boldsymbol s}}

\newcommand{\bxi}{{\boldsymbol \xi}}
\newcommand{\balpha}{{\boldsymbol \alpha}}
\newcommand{\bbeta}{{\boldsymbol \beta}}
\newcommand{\bgamma}{{\boldsymbol \gamma}}
\newcommand{\btheta}{{\boldsymbol \theta}}
\newcommand{\bTheta}{{\boldsymbol \Theta}}
\newcommand{\bvartheta}{{\boldsymbol \vartheta}}
\newcommand{\bepsilon}{{\boldsymbol \epsilon}}
\newcommand{\bvarepsilon}{{\boldsymbol \varepsilon}}
\newcommand{\bomega}{{\boldsymbol \omega}}
\newcommand{\bmu}{{\boldsymbol \mu}}
\newcommand{\bSigma}{{\boldsymbol \Sigma}}

\newcommand{\corr}{\mbox{Corr}}
\newcommand{\var}{{\mbox{Var}}}
\newcommand{\Var}{{\mbox{Var}}}
\newcommand{\cov}{{\mbox{Cov}}}
\newcommand{\diag}{{\mbox{diag}}}
\newcommand{\bUmis}{{\boldsymbol U}_{{\rm mis}}}
\newcommand{\bYmis}{{\boldsymbol Y}_{{\rm mis}}}
\newcommand{\bYobs}{{\boldsymbol Y}_{{\rm obs}}}
\newcommand{\tbXmis}{{\tilde{\boldsymbol X}}_{{\rm  mis}}}
\newcommand{\tbxmis}{{\tilde{\boldsymbol x}}_{{\rm  mis}}}
\newcommand{\bXmis}{{\boldsymbol X}_{{\rm mis}}}
\newcommand{\bXobs}{{\boldsymbol X}_{{\rm obs}}}
\newcommand{\bZmis}{{\boldsymbol Z}_{{\rm mis}}}
\newcommand{\bZobs}{{\boldsymbol Z}_{{\rm obs}}}
\newcommand{\tXmis}{{\tilde{X}}_{{\rm  mis}}}
\newcommand{\txmis}{{\tilde{x}}_{{\rm  mis}}}
\newcommand{\Xmis}{X_{{\rm  mis}}}
\newcommand{\Xobs}{X_{{\rm  obs}}}
\newcommand{\bymis}{{\boldsymbol y}_{{\rm  mis}}}
\newcommand{\byobs}{{\boldsymbol y}_{{\rm  obs}}}
\newcommand{\ymis}{y_{{\rm  mis}}}
\newcommand{\yobs}{y_{{\rm  obs}}}

\newcommand{\JRSSB}{Journal of the Royal Statistical Society: Series B}
\newcommand{\indep}{\rotatebox[origin=c]{90}{$\models$}}
\newcommand{\red}{\textcolor{red}}
\newcommand{\blue}{\textcolor{blue}}
\newcommand{\notetodo}[1]{\textcolor{blue}{[TODO: #1]}}

% Preserve the existing independent, section-based theorem counters.
\newtheorem{theorem}{Theorem}[section]
\newtheorem{lemma}{Lemma}[section]
\newtheorem{corollary}{Corollary}[section]
\newtheorem{formula}{Formula}[section]
\newtheorem{proposition}{Proposition}[section]
\theoremstyle{definition}
\newtheorem{definition}{Definition}[section]
\newtheorem{example}{Example}[section]
\newtheorem{assumption}{Assumption}
\theoremstyle{remark}
\newtheorem{remark}{Remark}

% Journal abbreviations used by the existing bibliography.
\def\AJE{American Journal of Epidemiology}
\def\ANNALS{Annals of Statistics}
\def\BIOK{Biometrika}
\def\BMCS{Biometrics}
\def\STATMED{Statistics in Medicine}
\def\STATSCI{Statistical Science}
\def\JSPI{Journal of Statistical Planning \& Inference}
\def\JRSSB{Journal of the Royal Statistical Society, Series B}
\def\COMMS{Communications in Statistics, Series A}
\def\JQT{Journal of Quality Technology}
\def\TECH{Technometrics}
\def\JASA{Journal of the American Statistical Association}
\def\CSDA{Computational Statistics \& Data Analysis}
\def\JCGS{Journal of Computational and Graphical Statistics}
\def\JCB{Journal of Computational Biology}
\def\BIOINF{Bioinformatics}
\def\JAMA{Journal of the American Medical Association}
\def\JNUTR{Journal of Nutrition}
\def\LETTERS{Letters in Probability and Statistics}
\def\JABES{Journal of Agricultural, Biological and Environmental Statistics}
\def\SCAN{Scandinavian Journal of Statistics}
\def\PNAS{Proceedings of the National Academy of Sciences USA}
\def\JMLR{Journal of Machine Learning Research}
\def\JSCS{Journal of Statistical Computation and Simulation}


\title{Note on the Proofs for NE-EFI}
\author{Wen-Hung Wang}
\date{\today}

\begin{document}
\maketitle
\tableofcontents
\bigskip

\section{Probability}
\subsection{Some Definitions}
\begin{definition}[Distribution Law]
    Given a probability space $(\Omega, \mathcal{A}, P)$, a \emph{distribution (law)} $\mu: \mathcal{B} \to [0,1]$ of an r.v. $X: \Omega \to \mathbb R$ is defined as 
\[
\mu(B) = P(\underbrace{X^{-1}(B)}_{\in \mathcal A}) = P(\{\omega \in \Omega: X(\omega) \in B\}) = \underbrace{P(X \in B)}_{\text{for convience}},
\]
which is a set function. Note that $\mu$ is a \emph{measure} ($\geq 0,~\mu(\varnothing) = 0$, countable additivity). For a absolutely continuous $\mu$, we may write
\begin{equation}
    \label{eq:abs_conti}
    \mu(B) = P(X \in B) = \int_B \textcolor{blue}{f(x)} \textcolor{red}{dx}
\end{equation}
where \textcolor{blue}{$f(x)$ is the density/Radon-Nikodym derivative} and \textcolor{red}{$dx$ is the Lebesgue measure} (counting measure if $X$ discrete). We call 
\[
\mathcal{L}(X) = \textcolor{blue}{f(x)} \textcolor{red}{dx}
\]
the law of $X$, sometimes denoted as $\mathcal{L}(X) =P(dx) = d\mu= f(x)dx$. This is a handy notation to denote a probability measure.
\end{definition}

\begin{definition}[Conditional Law]
    $\mathcal{L}(X|A) =  f(x|A) dx$ is a \emph{conditional law} if it satisfies for all $t$,
    \[
    P(X \geq t | A) = \int_{t}^\infty \underbrace{\mathcal{L}(X|A)}_{\text{prob. measure}} = \int_{t}^\infty f(x|A) dx
    \]
\end{definition}

\begin{definition}[Pushforward]
    Define measurable spaces $(\Omega_1, \mathcal{A}_1)$ and $(\Omega_2, \mathcal{A}_2)$. Let $\mu: \mathcal{A}_1 \to [0, +\infty]$ and $f: \Omega_1 \to \Omega_2$ be a measurable function. Then the \emph{pushforward measure} $f_{\#}(\mu): \mathcal{A}_2 \to [0, +\infty]$ can be defined as, for $A \in \mathcal{A}_2$
    \[
    f_{\#}(\mu)(A) = \mu(f^{-1}(A)). 
    \]
    Note that measurability of $f$ implies that $f^{-1}(A) \in \mathcal{A}_1$.    
\end{definition}

\begin{theorem}[Change of Variable for Pushforward]
    A measurable function $g$ on $\Omega_2$ is integrable wrt $f_{\#}(\mu)$ iff the following holds
    \[
    \int_{\Omega_2} g~d(f_{\#}(\mu)) = \int_{\Omega_1} g \circ f~d\mu
    \]
    where $\Omega_1 = f^{-1}(\Omega_2)$.
    
\end{theorem}

\begin{remark}[EFI as Pushforward]
    EFD in (15) in \cite{LiangKS2024EFI} is in essence a pushforward measure with $f=G$ and $\mu = dP^*_n(z|\bY_n,\bX_n)$.
\end{remark}

% \begin{lemma}[Scheffe's Lemma; \cite{scheffe1947useful}]
% \label{lemma:Scheffe}
% If a sequence of densities $\{\pi_n(z)\}_{n=1}^{\infty}$ satisfies $\pi_n(z)\to \pi(z)$ for almost every $z\in\mathbb{R}$, then a sufficient condition for
% \[
% \int_B \pi_n(z)\,dz \;\longrightarrow\; \int_B \pi(z)\,dz
% \]
% to hold uniformly over all Borel sets $B\subset\mathbb{R}$ is that $\pi$ itself be a (probability) density on $\mathbb{R}$.
%   %  If for a sequence of density $\{\pi_n(z)\}_{n=1}^{\infty}$, we have $\lim_{n \to \infty}\pi_n(z) = \pi(z)$ for almost every $z \in \mathbb{R}$, then a sufficient condition that $\lim_{n \to \infty}\int_B \pi_n(z)~dz = \int_B \pi(z)~dz$ holds uniformly for all Borel sets $B$ in $\mathbb{R}$ is that $\pi(z)$ be a density.
% \end{lemma}

% \begin{remark}
%     Lemma \ref{lemma:Scheffe} essentially says that a.s. convergence of density implies uniform convergence of CDF.
% \end{remark}

\subsection{$W_1$ Distance and Its Properties.}

\begin{definition}[1-Wasserstein Distance] Given two probability measures $P, Q$, the \emph{1-Wasserstein distance} is defined as
\[
W_1(P,Q) = \textcolor{blue}{\inf_{\gamma \in \Gamma(P,Q)}} E_{X,Y \sim \gamma} [d(X,Y)]
\]
where $d$ is the metric used in metric space, $\gamma$ is called a \emph{coupling}, and $\Gamma(P,Q)$ is the set of all joint distributions (transport plans) that have $P$ and $Q$ as their marginals. E.g., $P(dx) = \int_y \gamma(dx,dy)$.
\end{definition}

\begin{remark}
    {A useful technique to bound $W_1$ is to define $X,Y \sim \gamma^*$ where $\gamma^*$ is the optimal coupling}
\end{remark}

\paragraph{Properties.} Below are several properties of $W_1$:
\begin{enumerate}
    \item $W_1$ is smooth and differentiable, even if $P,Q$ not overlapping.
    \item Symmetric, $W_1(P,Q) = W_1(Q,P)$.
    \item \emph{Kantorovich-Rubinbstein Duality (KR duality)}:
    \[
    W_1(P,Q) = \textcolor{blue}{\sup_{|\|f\|_L \leq 1}}\left\{E_{X\sim P}[f(X)] - E_{Y\sim Q}[f(Y)] \right\}
    \]
    where the \textcolor{blue}{supremum is taken over all $f$ that are 1-Lipschitz continuous.} This implies
    \[
    |E_{X\sim P}[f(X)] - E_{Y\sim Q}[f(Y)]| \leq  \|f\|_LW_1(P,Q)
    \]
\end{enumerate}

\begin{remark}
    In practice, sometimes KR duality is more useful as once $f$ is established; the rest is easier than doing by definition.
\end{remark}

\paragraph{Proving $W_1$ or $W_2$?} In general, $W_1$ is weaker and often only requires bounded first moments. It is more suitable for heavy-tail distributions or general distribution distance. On the other hand, $W_2$ is stronger and considers second moments (geometry of distributions). Besides, $W_2$ penalizes large transport more:
\[
W_p(P,Q) = \inf_{\pi \in \Pi(P,Q)} E_{(X,Y) \sim \pi}[\|X-Y\|^p]^{1/p}.
\]
Assuming finite second moments, by Cauchy-Schwartz
\[
W_1(P, Q) \leq W_2(P, Q).
\]

\begin{lemma}[Convolution is non-expansive in \(\mathbb W_1\)]
\label{lem:conv_nonexp_mod}
For any probability measures \(Q,\widetilde Q\) on \(\mathbb R\),
\[
\mathbb W_1(K_h*Q,K_h*\widetilde Q)\le \mathbb W_1(Q,\widetilde Q).
\]
\end{lemma}

\begin{proof}
Let \((X,\widetilde X)\) be an optimal coupling of \(Q\) and \(\widetilde Q\), and let \(U\sim K_h\) be independent of both. Then
\(X+U\sim K_h*Q\), \(\widetilde X+U\sim K_h*\widetilde Q\), and
\[
E|X+U-(\widetilde X+U)|=E|X-\widetilde X|.
\]
Hence
\[
\mathbb W_1(K_h*Q,K_h*\widetilde Q)
\le
E|X+U-(\widetilde X+U)|
=
\mathbb W_1(Q,\widetilde Q).
\qedhere
\]
\end{proof}

\begin{lemma}[A truncation inequality for $\mathbb W_1$] \label{lem:W1_trunc_mod}
Let $X\sim\mu$ and $Y\sim\nu$ be any coupling with $E|X|+E|Y|<\infty$.
For an event $A$ with $P(A)>0$, write $\mu_A=\mathcal L(X\mid A)$ and $\nu_A=\mathcal L(Y\mid A)$.
Then
\[
{\mathbb{W}_1(\mu,\nu)\ \le\ \mathbb{W}_1(\mu_A,\nu_A)\ +\ E(|X| \mathbb{I}_{A^c})\ +\ E(|Y| \mathbb{I}_{A^c}).}
\]
where $\mathbb{I}_A$ denotes the indicator function on event $A$.
In particular, if $E X^2+E Y^2<\infty$ then
\[
{\mathbb{W}_1(\mu,\nu)\ \le\ \mathbb{W}_1(\mu_A,\nu_A)\ +\ \sqrt{E X^2}\sqrt{P(A^c)}\ +\ \sqrt{E Y^2}\sqrt{P(A^c)}.}
\]
\end{lemma}
\begin{proof}
Let $f$ be 1-Lipschitz with $f(0)=0$, so $|f(x)|\le |x|$.
Then
\[
|E f(X)-E f(Y)|
\le
|E(f(X)\mid A)-E(f(Y)\mid A)|P(A)
+ E(|X| \mathbb{I}_{A^c})+E(|Y| \mathbb{I}_{A^c}).
\]
Divide by $P(A)$ inside the first term to recognize $E(f(X)\mid A)=E_{\mu_A}f$, etc., apply $P(A) \leq 1$,
take the supremum over such $f$ (Kantorovich-Rubinstein duality), and use the Cauchy-Schwarz inequality for the second inequality.
\end{proof}

\begin{remark}
    {Lemma \ref{lem:W1_trunc_mod} is quite useful in bounding $W_1$ by defining an event $A$ with $A^c$ has small probability (tail argument). Then bound the first term.}
\end{remark}

\section{Jacobian and Gibbs Density Control in EFI}
This section provides review for some math tools. Some statements could be a bit imprecise.

\subsection{Jacobian}
\begin{theorem}[Change of Variable Formula]
\label{thm:change_var}
    Let $T$ be an one-to-one $C^1$ mapping of an open set $E \subset \mathbb{R}^k$ into $\mathbb{R}^k$ (same dimension) such that the Jacobian $J_T(\bx) = \det(\nabla_{\bx}T(\bx)) \neq {0}$. If $f$ is a continuous function supported on a compact set lying in $T(E)$, then
    \[
    \int_{\mathbb{R}^k}f(\by)d\by = \int_{\mathbb{R}^k} f(T(\bx)) |J_T(\bx)|d\bx
    \]
\end{theorem}

\begin{remark}
    For example, \emph{polar coordinate}.
\end{remark}

\paragraph{Gram Determinant.}
In Riemannian geometry, given an embedded $k$-dimensional Riemannian manifold $M \subset \mathbb{R}^n$ and a parametrization $\phi: U \rightarrow M$ for $\left(x_1, \ldots, x_k\right) \in U \subset \mathbb{R}^k$, the \emph{volume form} $\omega$ on $M$ induced by the embedding may be computed using the Gramian of the coordinate tangent vectors:
$$
\omega=\textcolor{red}{d\text{Vol}_M(\phi)} = \underbrace{\sqrt{\operatorname{det} G}}_{\approx \text{ Jacobian}} \textcolor{blue}{\underbrace{d x_1 \cdots d x_k}_{\text{ Lebesgue measure}}}, \quad G=\left[\left\langle\frac{\partial \phi}{\partial x_i}, \frac{\partial \phi}{\partial x_j}\right\rangle\right] .
$$
where $d\text{Vol}_M(\phi)$ is sometimes called \emph{intrinsic volume measure} and $x_i$'s are \emph{local coordinates}.

\begin{enumerate}
    \item The above formulation provides a change of variable formula from Lebesgue measure to intrinsic volume measure on a Riemannian manifold.

    \item $\phi$ maps local coordinates $(x_1,...,x_k)$ to a Riemannian manifold $M \subset \mathbb{R}^n$. \textcolor{red}{$\phi$ is $s$ in the NE-EFI paper or $m$ in the EFI paper given by the tubular neighborhood theorem.}

    \item One of the purposes using this formula is to \textcolor{red}{convert region of integral from chart domain to $M$.}
\end{enumerate}


% \begin{remark}[Change of Variable Used in Lemma D.14]
%     Let's revisit how we use these properties in the NE-EFI paper. We simplify the notations a bit. We note that $\bar P(z) = P(z | d(Z,\mathcal{Z}_n) \leq r)$ is a truncated probability measure on $\mathcal{N}_r(\mathcal{Z}_n)$. 
%     \[
%     \bar P(z) = \frac{P(Z, d(Z,\mathcal{Z}_n) \leq r)}{P( d(Z,\mathcal{Z}_n) \leq r)} := \frac{A}{B}
%     \]
%     Since we assume that $z \in \mathcal{N}_r(\mathcal{Z}_n)$ admits a parametrization $z = \Phi(s(u),t)$ with $s: \mathbb{R}^p \to \mathbb{R}^n$ maps the local coordinates to $\mathcal{Z}_n$, we obtain 
%     \[
%     A = \prod_i \pi(\Phi(s(u), t)_i) \exp\left(-\frac{U(\Phi(s(u), t))}{\epsilon}\right)dz
%     \]
%     and $B$ is 
%     \[
%     B= \int_{\mathbb{R}^{n}} \prod_i \pi(\Phi(s(u), t)_i) \exp\left(-\frac{U(\Phi(s(u), t))}{\epsilon}\right)dz.
%     \]
%     The normalizing constants are canceled out. Now, we need to change the variable $z$ to $(u,t)$ (not $(s(u),t)$). Applying the change of variable formula in Theorem \ref{thm:change_var} to $A$ and $B$ gives
%     \[
%      A = \prod_i \pi(\Phi(s(u), t)_i) \exp\left(-\frac{U(\Phi(s(u), t))}{\epsilon}\right) J(u,t) dudt
%     \]
%     and 
%     \[
%      B = \textcolor{blue}{\int_{u \in U}}\int_{t \in B}\prod_i \pi(\Phi(s(u), t)_i) \exp\left(-\frac{U(\Phi(s(u), t))}{\epsilon}\right) J(u,t) dt\textcolor{blue}{du}
%     \]
%     where $J(u,t) = |\det({D_{u,t}\Phi(s(u),t)})| = \sqrt{\det({D_{u,t}\Phi(s(u),t)}^\top{D_{u,t}\Phi(s(u),t))}}$. Next, we want the region of integration being on the manifold $\mathcal{Z}_n$, meaning that we need to use change of variable on manifold $\mathcal{Z}_n$ and integrate wrt its intrinsic volume measure. We therefore have $d\text{Vol}_{\mathcal{Z}_n}(s) = \sqrt{\det(G)} du$ where $G = D_u s(u)^\top D_u s(u)$. As a result, 
%     \[
%     B = \textcolor{blue}{\int_{s \in \mathcal{Z}_n}}\int_{t \in B}\prod_i \pi(\Phi(s(u), t)_i) \exp\left(-\frac{U(\Phi(s(u), t))}{\epsilon}\right) \textcolor{red}{\frac{J(u,t)}{\sqrt{\det(G)}}} dt\textcolor{blue}{d\text{Vol}_{\mathcal{Z}_n}(s)}
%     \]
%     In short
%     \begin{enumerate}
%         \item Phase 1: Converting to a Lebesgue measure integration.
%         \item Phase 2: Leverage the relationship between intrinsic measure and Lebesgue measure to perform manifold change of variable.
%     \end{enumerate}
% \end{remark}

\begin{remark}[Change of Variable Used in Lemma D.14]
    Let's revisit how these geometric properties are used to compute the expected normal deviation. To simplify the notation, let the unnormalized base measure be $\mu(dz) = \prod_i \pi(z_i) \exp(-U_n(z)/\epsilon) dz$. The truncated probability measure $\bar{P}$ on the tubular neighborhood $\mathcal{N}_r(\mathcal{Z}_n)$ can be expressed in terms of its differential:
    \[
    \bar{P}(dz) = \frac{\mu(dz)}{\int_{\mathcal{N}_r(\mathcal{Z}_n)} \mu(dz)} := \frac{A(dz)}{B}.
    \]
    
    We assume that $z \in \mathcal{N}_r(\mathcal{Z}_n)$ admits a local parametrization $z = \Psi_n(u,t) := \Phi_n(s(u),t)$, where $u \in U \subset \mathbb{R}^p$ are the local coordinates of the manifold, $s: U \to \mathcal{Z}_n$ maps the coordinates to the manifold, and $t \in B_r^{n-p}$ represents the normal coordinates. 
    
    \textbf{Phase 1: Ambient coordinates to local tubular coordinates.}\\
    We first apply the standard multivariate change of variables from $z \in \mathbb{R}^n$ to $(u,t) \in \mathbb{R}^p \times \mathbb{R}^{n-p}$. The unnormalized differential $A(dz)$ becomes:
    \[
    A(u,t) = \prod_i \pi(\Psi_n(u,t)_i) \exp\left(-\frac{U_n(\Psi_n(u,t))}{\epsilon}\right) \textcolor{blue}{J_{\Psi_n}(u,t) \,du \,dt},
    \]
    where $J_{\Psi_n}(u,t) = |\det(D_{(u,t)}\Psi_n(u,t))|$ is the absolute value of the determinant of the $n \times n$ Jacobian matrix of $\Psi_n$. 
    
    Consequently, the normalizing constant $B$ over the tubular neighborhood is:
    \[
    B = \textcolor{blue}{\int_{U}}\int_{B_r^{n-p}} \prod_i \pi(\Psi_n(u,t)_i) \exp\left(-\frac{U_n(\Psi_n(u,t))}{\epsilon}\right) J_{\Psi_n}(u,t) \textcolor{blue}{\,dt \,du}.
    \]

    \textbf{Phase 2: Local coordinates to intrinsic manifold integration.}\\
    Next, we want to integrate over the manifold $\mathcal{Z}_n$ with respect to its intrinsic volume measure, $d\mathrm{Vol}_{\mathcal{Z}_n}(s)$, rather than the parameter space $U$. The volume form on the submanifold is given by $d\mathrm{Vol}_{\mathcal{Z}_n}(s) = \sqrt{\det(G(u))} \,du$, where $G(u) = D_u s(u)^\top D_u s(u)$ is the metric tensor induced by the embedding. 
    
    Rearranging this gives $du = \textcolor{red}{\frac{1}{\sqrt{\det(G(u))}}} \textcolor{blue}{d\mathrm{Vol}_{\mathcal{Z}_n}(s)}$. Substituting this into our integral for $B$ yields:
    \[
    B = \textcolor{blue}{\int_{\mathcal{Z}_n}}\int_{B_r^{n-p}}\prod_i \pi(\Phi_n(s, t)_i) \exp\left(-\frac{U_n(\Phi_n(s, t))}{\epsilon}\right) \textcolor{red}{\left[ \frac{J_{\Psi_n}(u,t)}{\sqrt{\det(G(u))}} \right]} dt \textcolor{blue}{\,d\mathrm{Vol}_{\mathcal{Z}_n}(s)}.
    \]
    
    The bracketed term $\textcolor{red}{\frac{J_{\Psi_n}(u,t)}{\sqrt{\det(G(u))}}}$ is precisely the definition of $J_n(s,t)$, the Jacobian factor that couples the normal integration with the intrinsic manifold integration. This can be seen as a geometric volume distortion factor (in this way, one can bypass the transformation to $U$.)
\end{remark}


\section{Proof of Stopping Rule}

\clearpage
\bibliography{ref}
\bibliographystyle{abbrv}
\end{document}
